\section{Importancia del estudio}
\subsection{Resultados esperados}
El artefacto principal de esta investigación es un \textbf{prototipo de sistema informático}: un asistente educativo bilingüe español--Quechua Collao ejecutable sin conexión a internet en dispositivos de bajo o mediano costo. Este prototipo no constituye un modelo de lenguaje nuevo ni realiza entrenamiento del modelo base desde cero, sino que implementa un \textbf{enfoque híbrido} que integra un modelo compacto pre-entrenado y cuantizado con un \textit{pipeline} RAG local disciplinado, recursos lingüísticos explícitos (glosarios, diccionarios, reglas de normalización) y una adaptación ligera opcional (QLoRA), adaptado específicamente para responder consultas pedagógicas en contextos EIB del sur andino peruano. Los productos derivados de su construcción y evaluación son:
\begin{itemize}
	\item Un prototipo funcional de asistente educativo bilingüe español--Quechua Collao, ejecutable sin conexión a internet en dispositivos de bajo o mediano costo, basado en un modelo compacto cuantizado con \textit{pipeline} RAG disciplinado y adaptación ligera.
	\item Una base local de contenidos curriculares, glosarios terminológicos, diccionarios y materiales EIB en español y Quechua Collao, organizados para su recuperación híbrida (léxica + densa) y orientados al nivel primaria (3.° a 6.° grado), con posibilidad de adaptación posterior a secundaria básica.
	\item Un corpus curado de recursos lingüísticos español--Quechua Collao, complementado con datos sintéticos y validado por especialistas hablantes nativos.
	\item Un reporte comparativo de las configuraciones experimentales (\textit{baseline} sin RAG, RAG disciplinado y enfoque híbrido completo) que documente la contribución de cada componente del sistema.
	\item Un reporte de evaluación técnica que documente métricas de latencia, consumo de memoria, tokens por segundo y estabilidad \textit{offline} del prototipo en hardware de bajo o mediano costo.
	\item Un reporte de evaluación lingüística y pedagógica con resultados de métricas automáticas (chrF++, RAGAS, Recall@k) y valoraciones de docentes y especialistas EIB sobre la utilidad educativa, pertinencia intercultural y adecuación pedagógica de las respuestas generadas.
\end{itemize}

\subsection{Contribuciones}
\begin{itemize}
	\item \textbf{Contribución tecnológica:} Arquitectura híbrida funcional de un asistente educativo \textit{offline} que integra detección de idioma, enrutamiento lingüístico, RAG local disciplinado, recursos lingüísticos explícitos, adaptación ligera (QLoRA) y generación de respuestas mediante un modelo compacto cuantizado a 4 bits, ejecutable en dispositivos de bajo o mediano costo sin conexión a internet.
	\item \textbf{Contribución educativa:} Base de conocimiento local organizada con materiales curriculares EIB en español y Quechua Collao, orientada a apoyar consultas pedagógicas de estudiantes y docentes en zonas rurales, junto con un corpus curado de recursos lingüísticos para la variedad Quechua Collao aplicados a tareas de educación intercultural bilingüe.
	\item \textbf{Contribución metodológica:} Diseño de comparación escalonada que cuantifica la contribución de cada componente del enfoque híbrido, combinando métricas automáticas de calidad lingüística (chrF++), calidad de recuperación (RAGAS) y criterios pedagógicos e interculturales validados por especialistas EIB, aplicable a otros proyectos de tecnología educativa en lenguas indígenas de bajos recursos del sur andino.
\end{itemize}

\subsection{Impacto social}
El desarrollo de un asistente educativo bilingüe offline tiene el potencial de beneficiar a docentes y estudiantes de comunidades quechuahablantes del sur andino peruano que actualmente carecen de herramientas tecnológicas adaptadas a su realidad lingüística y territorial. Al operar sin conexión a internet y en Quechua Collao, el prototipo podría contribuir a reducir la brecha digital educativa en zonas rurales, fortalecer el uso de la lengua materna en contextos de aprendizaje formal y apoyar los objetivos del modelo de servicio educativo intercultural bilingüe del Ministerio de Educación.

\section{Cronograma de actividades}
El desarrollo de la tesis se planifica en un horizonte de doce meses a partir de la aprobación del plan. Las actividades se organizan en seis fases según la metodología Design Science Research.

\begin{table}[H]
\centering
\caption{Cronograma de actividades por fases y meses}
\label{tab:cronograma}
\small
\begin{tabular}{@{}p{5.5cm}cccccccccccc@{}}
\toprule
\textbf{Actividad} & \textbf{1} & \textbf{2} & \textbf{3} & \textbf{4} & \textbf{5} & \textbf{6} & \textbf{7} & \textbf{8} & \textbf{9} & \textbf{10} & \textbf{11} & \textbf{12} \\
\midrule
\textbf{Fase 1: Preparación y permisos} & & & & & & & & & & & & \\
Revisión de literatura y estado del arte & $\bullet$ & $\bullet$ & & & & & & & & & & \\
Caracterización de requerimientos EIB & $\bullet$ & $\bullet$ & & & & & & & & & & \\
Permisos institucionales y validación ética & $\bullet$ & $\bullet$ & $\bullet$ & & & & & & & & & \\
\midrule
\textbf{Fase 2: Diseño de la solución} & & & & & & & & & & & & \\
Diseño de la arquitectura híbrida \textit{offline} & & $\bullet$ & $\bullet$ & & & & & & & & & \\
Selección del modelo compacto base & & & $\bullet$ & $\bullet$ & & & & & & & & \\
\midrule
\textbf{Fase 3: Recursos y corpus} & & & & & & & & & & & & \\
Curación, depuración y normalización de corpus & & & $\bullet$ & $\bullet$ & $\bullet$ & & & & & & & \\
Glosarios, diccionarios y base curricular & & & & $\bullet$ & $\bullet$ & & & & & & & \\
Construcción del banco de consultas & & & & $\bullet$ & $\bullet$ & & & & & & & \\
\midrule
\textbf{Fase 4: Implementación} & & & & & & & & & & & & \\
Cuantización y configuración local del modelo & & & & & $\bullet$ & $\bullet$ & $\bullet$ & & & & & \\
Módulo de detección e integración de RAG local & & & & & & $\bullet$ & $\bullet$ & $\bullet$ & & & & \\
Adaptación ligera QLoRA (condicionada)* & & & & & & & & $\bullet$ & $\bullet$ & & & \\
Integración final del prototipo & & & & & & & & $\bullet$ & $\bullet$ & & & \\
\midrule
\textbf{Fase 5: Evaluación} & & & & & & & & & & & & \\
Diseño de rúbricas y experimentos E1--E7 & & & & & & & & & $\bullet$ & & & \\
Evaluación técnica y lingüística & & & & & & & & & $\bullet$ & $\bullet$ & & \\
Prueba piloto exploratoria y juicio de expertos EIB & & & & & & & & & & $\bullet$ & $\bullet$ & \\
Análisis y procesamiento de resultados & & & & & & & & & & & $\bullet$ & \\
\midrule
\textbf{Fase 6: Documentación} & & & & & & & & & & & & \\
Redacción del informe de tesis & & & & & & & & & & $\bullet$ & $\bullet$ & $\bullet$ \\
Revisión final con asesor y sustentación & & & & & & & & & & & & $\bullet$ \\
\bottomrule
\end{tabular}
\end{table}

\textit{*Nota: La adaptación QLoRA (Fase 4) se considera opcional y condicionada a la disponibilidad de recursos de cómputo y a la escala del corpus curricular validado.}


\section{Presupuesto}
\label{sec:presupuesto}
El presupuesto estimado para el desarrollo de la tesis se detalla a continuación. Se consideran los recursos necesarios para la implementación del prototipo, la evaluación y la documentación.

\begin{table}[H]
\centering
\caption{Presupuesto estimado del proyecto}
\label{tab:presupuesto}
\begin{tabular}{@{}p{7cm}rr@{}}
\toprule
\textbf{Rubro} & \textbf{Cantidad} & \textbf{Costo (S/.)} \\
\midrule
\multicolumn{3}{@{}l}{\textbf{Recursos tecnológicos}} \\
Dispositivo de prueba (laptop de gama media) & 1 & 2\,579.00 \\
Almacenamiento externo (SSD 1 TB) & 1 & 500.00 \\
Servicio de cómputo en la nube (cuantización y vectorización) & Puntual & 200.00 \\
Adaptación QLoRA (GPU por demanda) & Puntual & 150.00 \\
\midrule
\multicolumn{3}{@{}l}{\textbf{Recursos académicos}} \\
Material bibliográfico y acceso a bases de datos & Global & 300.00 \\
Impresiones y empastados & Global & 200.00 \\
\midrule
\multicolumn{3}{@{}l}{\textbf{Trabajo de campo}} \\
Transporte a zonas rurales (validación) & 5 visitas & 750.00 \\
Viáticos para trabajo de campo & 5 días & 500.00 \\
Material para talleres con docentes EIB & Global & 200.00 \\
\midrule
\multicolumn{3}{@{}l}{\textbf{Servicios profesionales}} \\
Consulta con especialistas en Quechua Collao & 3 sesiones & 450.00 \\
Validación lingüística por hablantes nativos & 5 sesiones & 500.00 \\
\midrule
\multicolumn{3}{@{}l}{\textbf{Otros}} \\
Imprevistos (10\%) & Global & 633.00 \\
\midrule
\textbf{TOTAL} & & \textbf{6\,962.00} \\
\bottomrule
\end{tabular}
\end{table}

\textit{Nota:} El uso de servicios de cómputo en la nube se restringe exclusivamente a tareas puntuales de desarrollo experimental (cuantización y afinamiento ligero) y no forma parte del funcionamiento final del asistente, el cual opera de forma 100\% offline. Los honorarios de especialistas y hablantes nativos corresponden a retribuciones por sesiones de validación lingüística y pedagógica del corpus y del prototipo.

